\documentclass{article}
\usepackage{amsmath,amssymb,amsthm}
\usepackage{algorithm}
\usepackage[noend]{algpseudocode}
\usepackage{enumitem}
\usepackage{hyperref}
\newtheorem{theorem}{Theorem}
\newtheorem{lemma}{Lemma}
\newtheorem{assumption}{Assumption}
\newtheorem{corollary}{Corollary}
\newcommand{\R}{\mathbb{R}}
\newcommand{\E}{\mathbb{E}}

\begin{document}

\section*{Heavy–Ball Momentum Method}

\section*{Mathematical Formulation}
Let $f:\R^{d}\to\R$ be differentiable.  Given step size $\alpha>0$ and momentum parameter $\beta\in[0,1)$,
the Heavy–Ball iteration reads
\[
\boxed{
\begin{aligned}
x_{k+1}&=x_{k}-\alpha \nabla f(x_{k})+\beta\,(x_{k}-x_{k-1}),\qquad k\ge 0,\\
x_{-1}&=x_{0}\;\;(\text{or any fixed vector}).
\end{aligned}}
\tag{HB}
\]
The momentum term uses the previous \emph{position difference} $x_{k}-x_{k-1}$ rather than a look‑ahead gradient.

\section*{Pseudocode}
\begin{algorithm}[H]
\caption{Heavy–Ball Momentum}
\begin{algorithmic}[1]
\Require Initial point $x_{0}\in\R^{d}$, step size $\alpha>0$, momentum $\beta\in[0,1)$,
maximum iterations $K$.
\State $x_{-1}\gets x_{0}$ \Comment{initial previous iterate}
\For{$k=0$ \textbf{to} $K-1$}
    \State $g_{k}\gets \nabla f(x_{k})$
    \State $x_{k+1}\gets x_{k}-\alpha g_{k}+\beta\,(x_{k}-x_{k-1})$
\EndFor
\State \Return $x_{K}$
\end{algorithmic}
\end{algorithm}

\section*{Dry Run Example}
Consider the univariate quadratic
\[
f(w)=(w-3)^{2},\qquad \nabla f(w)=2(w-3).
\]
Choose $\alpha=0.1$, $\beta=0.5$, and $x_{0}=0$.  Initialise $x_{-1}=x_{0}=0$.

\begin{tabular}{c|c|c|c}
Iteration $k$ & $g_{k}= \nabla f(x_{k})$ & $x_{k+1}=x_{k}-\alpha g_{k}+\beta(x_{k}-x_{k-1})$ & $f(x_{k+1})$\\\hline
0 & $2(0-3)=-6$ & $0-0.1(-6)+0.5(0-0)=0.6$ & $(0.6-3)^{2}=5.76$\\
1 & $2(0.6-3)=-4.8$ & $0.6-0.1(-4.8)+0.5(0.6-0)=0.6+0.48+0.3=1.38$ & $(1.38-3)^{2}=2.6244$\\
2 & $2(1.38-3)=-3.24$ & $1.38-0.1(-3.24)+0.5(1.38-0.6)=1.38+0.324+0.39=2.094$ & $(2.094-3)^{2}=0.8209$
\end{tabular}

The objective value decreases rapidly, illustrating the acceleration effect of the momentum term.

\newpage
\section*{Assumptions}
\begin{assumption}[Strong Convexity]\label{ass:sc}
$f$ is $\mu$‑strongly convex, i.e.
\[
f(y)\ge f(x)+\langle\nabla f(x),y-x\rangle+\frac{\mu}{2}\|y-x\|^{2},
\qquad\forall x,y\in\R^{d},
\]
with $\mu>0$.
\end{assumption}
\begin{assumption}[Smoothness]\label{ass:smooth}
$f$ has $L$‑Lipschitz continuous gradient:
\[
\|\nabla f(x)-\nabla f(y)\|\le L\|x-y\|,
\qquad\forall x,y\in\R^{d},
\]
with $L\ge\mu$.
\end{assumption}
\begin{assumption}[Parameter Choice]\label{ass:param}
The step size $\alpha$ and momentum $\beta$ satisfy
\[
0<\alpha<\frac{4}{(\sqrt{L}+\sqrt{\mu})^{2}},\qquad
0\le\beta<\bigl(1-\sqrt{\alpha\mu}\,\bigr)^{2}.
\tag{A3}
\]
\end{assumption}

\section*{Main Convergence Theorem}
\begin{theorem}[Linear Convergence of Heavy–Ball]\label{thm:linear}
Under Assumptions \ref{ass:sc}–\ref{ass:param}, let $x^{\star}=\arg\min f$.  
Define the Lyapunov vector
\[
z_{k}\;:=\;\begin{bmatrix}x_{k}-x^{\star}\\ x_{k-1}-x^{\star}\end{bmatrix}\in\R^{2d}.
\]
Then there exists a constant $\rho\in(0,1)$ depending only on $\alpha,\beta,\mu,L$ such that
\[
\|z_{k}\|\le\rho^{\,k}\,\|z_{0}\|,\qquad\forall k\ge0,
\tag{3}
\]
and consequently
\[
\|x_{k}-x^{\star}\|\le\rho^{\,k}\,\|x_{0}-x^{\star}\| .
\tag{4}
\]
In particular, for the optimal choice
\[
\alpha^{\star}= \frac{4}{(\sqrt{L}+\sqrt{\mu})^{2}},\qquad
\beta^{\star}= \Bigl(\frac{\sqrt{L}-\sqrt{\mu}}{\sqrt{L}+\sqrt{\mu}}\Bigr)^{2},
\]
the rate constant is
\[
\rho^{\star}= \frac{\sqrt{\kappa}-1}{\sqrt{\kappa}+1},\qquad
\kappa:=\frac{L}{\mu}.
\tag{5}
\]
\end{theorem}

\section*{Theoretical Properties}
\begin{itemize}[leftmargin=*]
\item \textbf{Stability.} Condition (A3) guarantees the spectral radius of the iteration matrix (see Lemma~\ref{lem:spectral}) is $<1$, which is necessary and sufficient for asymptotic stability of the linear system governing $z_{k}$.
\item \textbf{Hyper‑parameter Sensitivity.} The admissible region (A3) shrinks as the condition number $\kappa$ grows; choosing $\alpha$ too large or $\beta$ too close to $1$ may violate \eqref{eq:spectral-bound} and cause divergence.
\item \textbf{Comparison.} For gradient descent ($\beta=0$) the optimal linear rate is $1-\mu\alpha$, i.e. $\rho_{\text{GD}}=1-2/(\kappa+1)$. Heavy–Ball attains the accelerated rate $\rho^{\star}= (\sqrt{\kappa}-1)/(\sqrt{\kappa}+1)$, which is strictly smaller for $\kappa>1$.
\end{itemize}

\section*{Discussion}
The Heavy–Ball method can be interpreted as a discretisation of a second‑order differential equation with friction.  The proof below follows the classical energy‑method approach: a suitable quadratic Lyapunov function is shown to contract at each iteration.  The analysis is deterministic; no stochasticity is involved.

\newpage
\section*{Preliminaries (Definitions and Lemmas)}
\begin{lemma}[Descent Lemma]\label{lem:descent}
If $f$ is $L$‑smooth then for any $x,y\in\R^{d}$,
\[
f(y)\le f(x)+\langle\nabla f(x),y-x\rangle+\frac{L}{2}\|y-x\|^{2}.
\tag{L1}
\]
\end{lemma}
\begin{lemma}[Quadratic Upper/Lower Bounds]\label{lem:quad}
For a $\mu$‑strongly convex, $L$‑smooth function,
\[
\frac{\mu}{2}\|x-x^{\star}\|^{2}\le f(x)-f(x^{\star})\le\frac{L}{2}\|x-x^{\star}\|^{2},
\qquad\forall x\in\R^{d}.
\tag{L2}
\]
\end{lemma}
\begin{lemma}[Spectral Bound]\label{lem:spectral}
Define the $2d\times2d$ matrix
\[
M:=\begin{bmatrix}
(1+\beta)I-\alpha \nabla^{2}f(x^{\star}) & -\beta I\\[2mm]
I & 0
\end{bmatrix}.
\]
If $\alpha,\beta$ satisfy (A3) then the spectral radius $\rho(M)<1$.
\end{lemma}
\begin{proof}[Proof of Lemma~\ref{lem:spectral}]
Since $f$ is quadratic at $x^{\star}$ with Hessian $H:=\nabla^{2}f(x^{\star})$ satisfying $\mu I\preceq H\preceq L I$, the eigenvalues of $M$ are the roots of
\[
\lambda^{2}-(1+\beta-\alpha\lambda_{i})\lambda+\beta=0,\qquad \lambda_{i}\in[\mu,L].
\tag{S1}
\]
Solving the quadratic gives
\[
\lambda_{i}^{\pm}=\frac{1+\beta-\alpha\lambda_{i}\pm\sqrt{(1+\beta-\alpha\lambda_{i})^{2}-4\beta}}{2}.
\tag{S2}
\]
Under (A3) one checks that $|\,\lambda_{i}^{\pm}\,|<1$ for all $\lambda_{i}\in[\mu,L]$; the worst case occurs at the endpoints $\lambda_{i}=\mu$ or $L$, yielding the bound $\rho(M)\le\rho^{\star}$ defined in \eqref{5}.  ∎
\end{proof}

\section*{Main Proof (Step‑by‑Step Derivation)}
We prove Theorem~\ref{thm:linear}.  Throughout we set $e_{k}:=x_{k}-x^{\star}$.

\begin{enumerate}
\item[Step 1.]  Write the HB update in error form:
\[
e_{k+1}=e_{k}-\alpha \nabla f(x_{k})+\beta(e_{k}-e_{k-1}).
\tag{1}
\]

\item[Step 2.]  Use the first‑order optimality condition $\nabla f(x^{\star})=0$ and add–subtract $\alpha\nabla f(x^{\star})$:
\[
e_{k+1}=e_{k}+\beta(e_{k}-e_{k-1})-\alpha\bigl[\nabla f(x_{k})-\nabla f(x^{\star})\bigr].
\tag{2}
\]

\item[Step 3.]  By smoothness, $\nabla f(x_{k})-\nabla f(x^{\star}) = H_{k}\,e_{k}$ for some symmetric $H_{k}$ satisfying $\mu I\preceq H_{k}\preceq L I$.  Substituting,
\[
e_{k+1}= \bigl(I-\alpha H_{k}+\beta I\bigr)e_{k}-\beta e_{k-1}.
\tag{3}
\]

\item[Step 4.]  Stack the two consecutive errors into a single vector
\[
z_{k}:=
\begin{bmatrix}e_{k}\\ e_{k-1}\end{bmatrix},
\qquad
z_{k+1}=A_{k}\,z_{k},
\tag{4}
\]
with
\[
A_{k}:=
\begin{bmatrix}
(1+\beta)I-\alpha H_{k} & -\beta I\\[1mm]
I & 0
\end{bmatrix}.
\tag{5}
\]

\item[Step 5.]  Because $H_{k}$ varies with $k$, we compare $A_{k}$ with the constant matrix $M$ defined in Lemma~\ref{lem:spectral}.  From $\mu I\preceq H_{k}\preceq L I$ we have
\[
(1+\beta)I-\alpha L I\;\preceq\;(1+\beta)I-\alpha H_{k}\;\preceq\;(1+\beta)I-\alpha\mu I .
\tag{6}
\]

\item[Step 6.]  Define the quadratic Lyapunov function
\[
\Phi_{k}:=z_{k}^{\top}P\,z_{k},
\tag{7}
\]
where $P\succ0$ will be chosen later.  Our goal is to show $\Phi_{k+1}\le\rho^{2}\Phi_{k}$ for some $\rho\in(0,1)$.

\item[Step 7.]  Compute
\[
\Phi_{k+1}=z_{k}^{\top}A_{k}^{\top}PA_{k}\,z_{k}.
\tag{8}
\]
We require
\[
A_{k}^{\top}PA_{k}\;\preceq\;\rho^{2}P\quad\forall\,k.
\tag{9}
\]

\item[Step 8.]  Choose $P$ as the solution of the discrete Lyapunov inequality for the worst‑case matrix $M$:
\[
M^{\top}PM\;\preceq\;\rho^{2}P,
\tag{10}
\]
with $\rho$ equal to the spectral radius $\rho(M)$.  Because $A_{k}$ lies between the two extremal matrices obtained by replacing $H_{k}$ with $\mu I$ or $L I$, inequality \eqref{9} follows from the monotonicity of the map $H\mapsto A(H)^{\top}PA(H)$ (a direct calculation shows the map is affine in $H$ and the ordering \eqref{6} preserves the inequality).  Hence
\[
\Phi_{k+1}\le\rho^{2}\Phi_{k}.
\tag{11}
\]

\item[Step 9.]  Unfolding \eqref{11} yields
\[
\Phi_{k}\le\rho^{2k}\Phi_{0}.
\tag{12}
\]
Since $P\succ0$, there exist constants $c_{1},c_{2}>0$ such that
\[
c_{1}\|z_{k}\|^{2}\le\Phi_{k}\le c_{2}\|z_{k}\|^{2}.
\tag{13}
\]
Combining \eqref{12}–\eqref{13} gives
\[
\|z_{k}\|\le\sqrt{\frac{c_{2}}{c_{1}}}\,\rho^{k}\|z_{0}\|.
\tag{14}
\]

\item[Step 10.]  The first block of $z_{k}$ is $e_{k}=x_{k}-x^{\star}$, therefore
\[
\|x_{k}-x^{\star}\|\le C\,\rho^{k}\|x_{0}-x^{\star}\|,\qquad C:=\sqrt{c_{2}/c_{1}}.
\tag{15}
\]
This establishes the linear rate \eqref{3}–\eqref{4}.  The explicit expression for $\rho$ follows from Lemma~\ref{lem:spectral}; for the optimal parameters $\alpha^{\star},\beta^{\star}$ the eigenvalues \eqref{S2} reduce to $\pm\rho^{\star}$ with
\[
\rho^{\star}= \frac{\sqrt{\kappa}-1}{\sqrt{\kappa}+1},
\tag{16}
\]
which completes the proof. ∎
\end{enumerate}

\section*{Rate Derivation (With Constants)}
From Lemma~\ref{lem:spectral} we have
\[
\rho = \max_{\lambda\in[\mu,L]}\Bigl|\frac{1+\beta-\alpha\lambda\pm\sqrt{(1+\beta-\alpha\lambda)^{2}-4\beta}}{2}\Bigr|.
\tag{R1}
\]
Choosing $\alpha$ and $\beta$ as in \eqref{5} gives
\[
\rho = \frac{\sqrt{L}-\sqrt{\mu}}{\sqrt{L}+\sqrt{\mu}} = \frac{\sqrt{\kappa}-1}{\sqrt{\kappa}+1}.
\tag{R2}
\]
Hence for all $k\ge0$,
\[
\|x_{k}-x^{\star}\|\le\Bigl(\frac{\sqrt{\kappa}-1}{\sqrt{\kappa}+1}\Bigr)^{k}\|x_{0}-x^{\star}\|.
\tag{R3}
\]
Equivalently, in terms of function values using Lemma~\ref{lem:quad},
\[
f(x_{k})-f(x^{\star})\le \frac{L}{2}\Bigl(\frac{\sqrt{\kappa}-1}{\sqrt{\kappa}+1}\Bigr)^{2k}\|x_{0}-x^{\star}\|^{2}.
\tag{R4}
\]

\section*{Special Cases}
\begin{itemize}[leftmargin=*]
\item \textbf{Exact Quadratic.} If $f$ is exactly quadratic, $H_{k}=H$ is constant and the iteration is linear; the bound \eqref{R3} is tight.
\item \textbf{Vanishing Momentum.} Setting $\beta=0$ reduces \eqref{HB} to gradient descent; the rate becomes $\rho=1-\alpha\mu$, recovered from \eqref{R1} with $\beta=0$.
\item \textbf{Over‑damped Regime.} If $\beta$ exceeds the bound in (A3) the discriminant in \eqref{S2} becomes negative and the spectral radius exceeds $1$, leading to divergence.  This illustrates the necessity of (A3).
\end{itemize}

\end{document}